\documentclass[12pt,a4paper]{article}
\usepackage[margin=1in]{geometry}
\usepackage{amsmath,amssymb}
\usepackage{hyperref}

\title{\textbf{The Cascade Series}}
\author{[Author]}
\date{March 2026}

\begin{document}
\maketitle
\thispagestyle{empty}

\section*{The Thought Experiment}

A black hole has a horizon. An infalling observer sees infinite time dilation there.
But the black hole evaporates in finite time. These two facts are in tension.

The resolution is asymptotic: there is no sharp horizon. Infalling matter never
crosses --- from outside, it asymptotically approaches the horizon; from the infaller's
perspective, the black hole evaporates and the horizon recedes. Both observers agree:
the infaller asymptotically tracks a shrinking shell where the rate of infall balances
the rate of evaporation.

This shell is the physical content of the ``horizon.'' On it sits a two-dimensional
spatial surface ($S^2$ for a Schwarzschild hole), and the time direction is the asymptotic
balance itself --- the direction along which infall and shrinkage compete. What if
this shell is a $2 + 1$ dimensional universe?

If a 4D black hole produces a $2 + 1$ shell universe, does a 5D black hole produce
a $3 + 1$ shell universe? Is that us? Are we living on the $S^3$ horizon of a 5D black
hole? If so, time is the forced direction --- the compactification that the cascade
must undergo, not a coordinate we choose. Moving orthogonally, along the shell,
slows progress in this forced direction but can never stop it completely. This is time
dilation, from the geometry of a shell.

If yes, then the 5D black hole sits inside a 6D space --- another shell of a 7D object.
Follow the tower upward: each shell is the horizon of the next. The tower terminates
at $d = \infty$: the infinite-dimensional unit ball, which has zero volume, zero surface
area, no interior. Geometric nothing. A natural starting point.

That was the end of the intuition. The question was: what happens if you work
back down from infinity?

\section*{The Hypothesis}

\begin{center}
\textbf{The infinite-dimensional unit ball, descended to four dimensions, is indistinguishable from our universe.}
\end{center}

It is tested by deriving, from the cascade's geometry alone, the cosmological constant, the dimension and signature of spacetime, quantum mechanics, and the
Standard Model gauge group and its symmetry breaking, three fermion generations,
fermion masses, and gauge coupling constants. It further derives the background
cosmological parameters --- including a Hubble constant between the two competing
measurements and an account of the DESI baryon acoustic oscillation observations
without dynamical dark energy. Every prediction is a test of the hypothesis.

\section*{The Series}

\textbf{Part I: Scale Variance from Orthogonality.} Pure mathematics. How much scale
variance can the unit ball produce from nothing? The slicing recurrence contains
one constant ($\sqrt{\pi}$), whose natural zero generates two distinguished dimensions at
$d = 19$ and $d = 217$, and a cascade invariant $I = 1.0996 \times 10^{-120}$ from the Gamma
function alone. The observed ratio $\rho_\Lambda/M_{\text{Pl}}^4 \approx 1.1 \times 10^{-120}$ was not the target. It
was the answer.

\textbf{Part II: Quantum Mechanics from the Cascade.} Does quantum mechanics
emerge naturally from the geometry of the descending unit ball? Consecutive slicing
axes are forced orthogonal, producing complex structure, Hilbert space, the Born
rule, and unitary evolution --- without quantum postulates.

\textbf{Part III: General Relativity from the Cascade.} Does general relativity emerge?
The cascade's complex spinor structure intersected with Lovelock uniqueness forces
$d = 4$, Lorentzian signature, and Einstein's equation. The cosmological constant
is $\Lambda = I$, the dark energy equation of state $w = -1$ exactly, both without free
parameters.

\textbf{Part IV: The Standard Model from the Cascade.} Does the Standard Model
emerge? Bott periodicity, Adams' theorem, and the hairy ball theorem jointly force
$\text{SU}(3) \times \text{SU}(2) \times \text{U}(1)$, its symmetry breaking pattern, three fermion generations,
and fourteen precision predictions --- all sub-3\%, zero free parameters.

\textbf{Part V: Cosmology from the Cascade.} Do the observed cosmological parameters
emerge? The cascade derives $\Omega_m$, $\Omega_b$, $H_0$, and $r_d$ from first principles. $H_0 =
70.65$~km/s/Mpc sits between the two competing measurements. The cascade's
$r_d = 142.3$~Mpc fits DESI DR2 BAO with $w = -1$ throughout; the apparent DESI
preference for $w \neq -1$ is a zero-parameter consequence of imposing the wrong ruler.

\bigskip

The hypothesis survives all five tests. The predictions separate into two populations: descent-dependent quantities ($\alpha_s$, $v$, $\Omega_m$, mass ratios, $m_\tau$, $m_W$) carry uniformly negative deviations of $-0.13\%$ to $-2.20\%$; geometric quantities ($m_H/m_W$, $\sin^2\theta_W$, $\theta_C$, $\Omega_b$) carry positive deviations of $+0.56\%$ to $+2.76\%$. The clean sign separation identifies the subleading cascade descent as the source of the dominant systematic, not computed in this series.

\end{document}
